\section{Summation and the bulk correlation}\label{sec:conclusion}

The preceding comparison gives a convergent amplitude after the
leading power and logarithm have been removed.  The fixed-geometry
comparison then transfers that amplitude from center magnetization
to the thermodynamic correlation.

\begin{proof}[Proof of Theorem~\ref{thm:main}]
Put \(\ell=\log L\) and
\[
                 c_N=g_{L^N}-2\pi\ell N.
\]
By \eqref{obs:green-constant}, \(c_N\) converges.  The
trajectory estimate in Equation~\eqref{trj:unbalanced-conclusion}, used at
\(j'=N-O(\log N)\) in its valid interior interval, gives
\[
 t_{j'}=\frac{1}{2\ell j'}
            +O\left(\frac{\log j'}{(j')^2}\right)
        =\frac{1}{2\ell N}
            +O\left(\frac{\log N}{N^2}\right).
\]
Since \(c_N\) is bounded, substituting this expression and
\(g_{L^{N+1}}-g_{L^N}=2\pi\ell+c_{N+1}-c_N\) into
Equation~\eqref{obs:consecutive-formula} yields
\begin{equation}\label{eq:magnetization-increment}
 \begin{split}
 \log M_{L^{N+1}}-\log M_{L^N}
  ={}&-\frac{\ell}{8}+\frac{1}{16N}
       -\frac{c_{N+1}-c_N}{16\pi}\\
     &+\frac{c_{N+1}-c_N}{32\pi\ell N}+\rho_N,
                \qquad \sum_N|\rho_N|<\infty .
 \end{split}
\end{equation}
Here \(a=8\pi\) fixes the two numerical coefficients.  The errors
are bounded by a constant times
\[
 N^{-2+4\epsilon}+N^{-1-\epsilon}
                       +\frac{\log N}{N^2},
\]
which is summable for the chosen \(0<\epsilon<1/10\).  The term
involving the error in \(t_{j'}\) times \(c_{N+1}-c_N\) is included
in the last summand because \(c_N\) is bounded.

The unweighted differences of \(c_N\) telescope.  Their weighted
series also converges: for \(M\ge N_0\),
\[
 \sum_{N=N_0}^M\frac{c_{N+1}-c_N}{N}
 =\frac{c_{M+1}}M-\frac{c_{N_0}}{N_0}
   +\sum_{N=N_0+1}^M c_N
                \left(\frac1{N-1}-\frac1N\right).
\]
The last sum converges absolutely since \(c_N\) is bounded, and the
first term tends to zero.  No bounded-variation assertion about
\((c_N)\) is needed.  Finally,
\(\sum_{N<N'}N^{-1}-\log N'\) has a finite limit.  Summing
Equation~\eqref{eq:magnetization-increment} proves that
\[
       \log M_{L^N}+\frac{\ell N}{8}-\frac{\log N}{16}
\]
has a finite real limit.  Therefore, for some \(b_M\in(0,\infty)\),
\begin{equation}\label{eq:center-asymptotic}
          M_{L^N}
     =b_M(L^N)^{-1/8}\bigl(\log L^N\bigr)^{1/16}(1+o(1)).
\end{equation}
The fixed change from \(N^{1/16}\) to
\((\log L^N)^{1/16}=(\ell N)^{1/16}\) is absorbed into \(b_M\).

This equivalent holds for every integer \(n\).  Given any sequence
\(n\to\infty\), let \(N=\lfloor\log_L n\rfloor\).  Every subsequence
has a further subsequence on which \(n/L^N\to d\in[1,L]\).
The center ratio \eqref{geo:center-ratio} in
Proposition~\ref{geo:comparison} gives
\(M_n/M_{L^N}\to d^{-1/8}\), while
\[
 \frac{n^{-1/8}(\log n)^{1/16}}
      {(L^N)^{-1/8}(\log L^N)^{1/16}}\longrightarrow d^{-1/8}.
\]
Equation~\eqref{eq:center-asymptotic} therefore holds along that
further subsequence.  The subsequence criterion proves it along
all integers.

The bulk-to-center limit \eqref{geo:bulk-ratio} in
Proposition~\ref{geo:comparison} now gives
\[
 C_{b_c}(r)
   =c_{\rm shape}M_r^2(1+o(1))
   =B_{\rm XY}\,r^{-1/4}(\log r)^{1/8}(1+o(1)),
 \qquad B_{\rm XY}=c_{\rm shape}b_M^2\in(0,\infty).
\]
The geometric limit concerns the thermodynamic correlation already
defined in Equation~\eqref{eq:critical-definition}; all cutoffs and
finite-square computations used only \(M_n\).  Thus the order of
limits in the theorem is the prescribed one.
\end{proof}
